% ═══════════════════════════════════════════════════════════════
% AB-07a: Comida (Lebensmittel)
% Abgeleitet aus LE-07a (Score 9.3/10)
% ═══════════════════════════════════════════════════════════════

\documentclass[11pt, a4paper]{article}

% ═══════════════════════════════════════════════════════════════
% SW-MODUS TOGGLE
% ═══════════════════════════════════════════════════════════════
% Setze \swmodustrue für Schwarz-Weiß-Druck
\newif\ifswmodus
\swmodusfalse    % Standard: Farbe

\usepackage[utf8]{inputenc}
\usepackage[T1]{fontenc}
\usepackage[ngerman]{babel}
\usepackage{helvet}
\renewcommand{\familydefault}{\sfdefault}

\usepackage[margin=2cm]{geometry}
\usepackage{enumitem}
\usepackage{tabularx}
\usepackage{booktabs}
\usepackage{multicol}

\usepackage{xcolor}
\usepackage{tcolorbox}
\tcbuselibrary{skins}

\usepackage{fancyhdr}
\usepackage{lastpage}
\usepackage{needspace}

% Farben - Basis
\definecolor{spanisch-color}{HTML}{c41e3a}
\definecolor{grau-color}{HTML}{888888}
\definecolor{hellgrau-color}{HTML}{f5f5f5}
\definecolor{orange-color}{HTML}{b35c00}

% SW-Farben (druckerfreundlich: keine Füllflächen, nur Linien)
\definecolor{sw-dark}{HTML}{000000}
\definecolor{sw-gray}{HTML}{666666}
\definecolor{sw-white}{HTML}{ffffff}

% Bedingte Farbzuweisung
\ifswmodus
    \colorlet{spanisch}{sw-dark}
    \colorlet{grau}{sw-gray}
    \colorlet{hellgrau}{sw-white}
    \colorlet{orange}{sw-dark}
\else
    \colorlet{spanisch}{spanisch-color}
    \colorlet{grau}{grau-color}
    \colorlet{hellgrau}{hellgrau-color}
    \colorlet{orange}{orange-color}
\fi

\newcommand{\fachfarbe}{spanisch}

% Variablen
\newcommand{\thema}{Comida -- Lebensmittel}
\newcommand{\sequenz}{07}
\newcommand{\block}{a}
\newcommand{\fach}{Spanisch}
\newcommand{\klasse}{7}
\newcommand{\maxpunkte}{24}

% Kopf-/Fußzeile
\pagestyle{fancy}
\fancyhf{}
\renewcommand{\headrulewidth}{0.4pt}
\renewcommand{\footrulewidth}{0.4pt}
\lhead{\fach{} -- Klasse \klasse}
\chead{AB-\sequenz\block}
\rhead{}
\lfoot{}
\rfoot{Seite \thepage\ von \pageref{LastPage}}

% Befehle
\newcommand{\punktebox}[1]{\hfill\fbox{\small \_/#1}}
\newcommand{\luecke}[1]{\rule{#1}{0.4pt}}
\newcommand{\es}[1]{\textcolor{\fachfarbe}{\textbf{#1}}}

% Merkkasten (SW-kompatibel)
\newtcolorbox{merkkasten}[1][]{
    colback=hellgrau,
    colframe=\fachfarbe,
    colbacktitle=white,
    coltitle=\fachfarbe,
    boxrule=1pt,
    arc=0pt,
    left=8pt, right=8pt, top=8pt, bottom=8pt,
    fonttitle=\bfseries,
    attach boxed title to top left={yshift=-2mm, xshift=5mm},
    boxed title style={boxrule=0pt, colback=white},
    title=#1
}

% Hinweis (SW-kompatibel)
\newtcolorbox{hinweis}{
    colback=white,
    colframe=orange,
    colbacktitle=white,
    coltitle=orange,
    boxrule=1pt,
    arc=0pt,
    left=5pt, right=5pt, top=5pt, bottom=5pt,
    fonttitle=\bfseries,
    attach boxed title to top left={yshift=-2mm, xshift=5mm},
    boxed title style={boxrule=0pt, colback=white},
    title=Hinweis
}

% Aufgabe (kein Seitenumbruch innerhalb)
\newenvironment{aufgabe}[1]{%
    \needspace{5\baselineskip}%
    \nopagebreak[4]%
    \vspace{10pt}
    \noindent\textbf{\textcolor{\fachfarbe}{Aufgabe #1}}
    \nopagebreak[4]%
    \vspace{5pt}
    \nopagebreak[4]%
    \begin{list}{}{\leftmargin=0pt}
    \item[]
}{%
    \end{list}
}

% ═══════════════════════════════════════════════════════════════
\begin{document}

% Kopfbereich
\begin{center}
    {\Large\bfseries\textcolor{\fachfarbe}{Arbeitsblatt: \thema}}
\end{center}

\vspace{5pt}

\noindent
\begin{tabularx}{\textwidth}{|X|X|X|}
\hline
\textbf{Name:} \luecke{4cm} &
\textbf{Klasse:} \klasse &
\textbf{Datum:} \luecke{2cm} \\
\hline
\end{tabularx}

\vspace{15pt}

% ═══════════════════════════════════════════════════════════════
% MERKKASTEN: Lebensmittel-Kategorien ($\rightarrow$ FZ1, FZ2)
% ═══════════════════════════════════════════════════════════════

\begin{merkkasten}[Merkkasten: Lebensmittel (La comida)]

\begin{tabularx}{\textwidth}{|l|X|}
\hline
\textbf{La fruta} (Obst) & \luecke{10cm} \\[0.8em]
\hline
\textbf{La verdura} (Gemüse) & \luecke{10cm} \\[0.8em]
\hline
\textbf{La carne / El pescado} & \luecke{10cm} \\[0.8em]
\hline
\textbf{Otros} (Andere) & \luecke{10cm} \\[0.8em]
\hline
\end{tabularx}

\vspace{0.5em}
\textit{Übertrage die Vokabeln aus dem Lernpfad oder von der Tafel!}
\end{merkkasten}

\vspace{10pt}

% ═══════════════════════════════════════════════════════════════
% AUFGABE 1: Zuordnung Bild-Wort ($\rightarrow$ FZ1, AFB I)
% ═══════════════════════════════════════════════════════════════

\begin{aufgabe}{1}
Verbinde die deutschen Wörter mit dem spanischen Begriff. \punktebox{5}
\end{aufgabe}

\begin{center}
\begin{tabular}{rcl}
Apfel & \luecke{3cm} & el plátano \\[0.8em]
Banane & \luecke{3cm} & la lechuga \\[0.8em]
Tomate & \luecke{3cm} & la manzana \\[0.8em]
Salat & \luecke{3cm} & la leche \\[0.8em]
Milch & \luecke{3cm} & el tomate \\
\end{tabular}
\end{center}

\vspace{15pt}

% ═══════════════════════════════════════════════════════════════
% AUFGABE 2: Kategorisierung ($\rightarrow$ FZ2, AFB I/II)
% ═══════════════════════════════════════════════════════════════

\begin{aufgabe}{2}
Ordne die Lebensmittel in die richtige Kategorie ein. \punktebox{6}
\end{aufgabe}

\begin{hinweis}
\textbf{Wörter:} el pollo, la naranja, el queso, la zanahoria, el pan, el pescado, el plátano, la cebolla, la carne, la leche
\end{hinweis}

\vspace{0.5em}

\begin{tabularx}{\textwidth}{|X|X|}
\hline
\textbf{La fruta} & \textbf{La verdura} \\[2.5em]
\hline
\textbf{La carne / El pescado} & \textbf{Otros} \\[2.5em]
\hline
\end{tabularx}

\vspace{15pt}

% ═══════════════════════════════════════════════════════════════
% AUFGABE 3: Lückentext Mengenangaben ($\rightarrow$ FZ3, AFB II)
% ═══════════════════════════════════════════════════════════════

\begin{aufgabe}{3}
Setze die richtige Mengenangabe ein. \punktebox{5}
\end{aufgabe}

\begin{hinweis}
\textbf{Mengenangaben:} un kilo de | medio kilo de | un litro de | una botella de | un paquete de
\end{hinweis}

\vspace{0.5em}

\begin{enumerate}[label=\alph*)]
    \item Quiero \luecke{3.5cm} manzanas.
    \item Necesito \luecke{3.5cm} tomates.
    \item Dame \luecke{3.5cm} leche.
    \item Quiero \luecke{3.5cm} agua.
    \item Necesito \luecke{3.5cm} arroz.
\end{enumerate}

\vspace{15pt}

% ═══════════════════════════════════════════════════════════════
% MERKKASTEN: Mengenangaben
% ═══════════════════════════════════════════════════════════════

\begin{merkkasten}[Merkkasten: Mengenangaben (Cantidades)]

\begin{tabular}{ll}
\es{un kilo de...} & ein Kilo... \\
\es{medio kilo de...} & ein halbes Kilo... \\
\es{un litro de...} & ein Liter... \\
\es{una botella de...} & eine Flasche... \\
\es{un paquete de...} & ein Paket/eine Packung... \\
\end{tabular}

\vspace{0.5em}
\textbf{Wichtig:} Nach Mengenangaben steht immer \es{de}!
\end{merkkasten}

\vspace{15pt}

% ═══════════════════════════════════════════════════════════════
% AUFGABE 4: Einkaufsliste schreiben ($\rightarrow$ FZ4, AFB III)
% ═══════════════════════════════════════════════════════════════

\begin{aufgabe}{4}
Schreibe eine Einkaufsliste auf Spanisch (mindestens 5 Produkte mit Mengenangaben). \punktebox{8}
\end{aufgabe}

\begin{hinweis}
Verwende verschiedene Kategorien (Obst, Gemüse, Fleisch, andere) und Mengenangaben!
\end{hinweis}

\vspace{0.5em}

\noindent\textbf{Mi lista de compras:}

\begin{enumerate}
    \item \luecke{8cm}
    \item \luecke{8cm}
    \item \luecke{8cm}
    \item \luecke{8cm}
    \item \luecke{8cm}
    \item \luecke{8cm} \textit{(Bonus)}
\end{enumerate}

\vspace{20pt}
\hrule
\vspace{10pt}

\begin{flushright}
\textbf{Punkte:} \luecke{1cm} / \maxpunkte
\end{flushright}

\end{document}
